\section{Entire harmonic functions and growth at every radius}
\label{n:sec-growth}

We use the one control sequence selected in
Section~\ref{n:sec-matching-growth}. For its metric $g=g_z$, the exact
outward maps satisfy \eqref{n:eq-exact-lines} with positive multipliers
$\sigma_i(j)$ and independent low columns $v_i(j)$. First we lift these
columns to compatible smooth boundary values on nested balls. We then
convert the strict phase means into the pointwise degree-$k$ bound.

\subsection{Compatible traces and entire functions}

With the graph maps of Proposition~\ref{n:prop-exact-graphs}, let
\[
 w_i(j)=\iota_jv_i(j)+V_jv_i(j)\in L^2(\mu).
\]
Equation \eqref{n:eq-graph-identity} and \(A_j=M_j(V_{j+1})^{-1}\) turn
\eqref{n:eq-exact-lines} into
\[
 T_jw_i(j+1)=\sigma_i(j)^{-1}w_i(j).
\]
Set
\[
 c_i(J)=1,\qquad
 c_i(j)=\prod_{\ell=J}^{j-1}\sigma_i(\ell),\qquad
 g_i(j)=c_i(j)w_i(j).
\]
Then \(T_jg_i(j+1)=g_i(j)\) for every \(j\geq J\).
These traces initially lie in \(L^2\).  The extension from the next larger
ball is smooth at the inner sphere by Lemma~\ref{n:lem-inward-transfer};
the transfer identity therefore also shows that each \(g_i(j)\) is a smooth
trace.  On \(B_g(0,e^j)\), take its classical harmonic Dirichlet extension
\cite[Theorem~6.14, p.~107]{GilbargTrudinger}.  The
extension on \(B_g(0,e^{j+1})\) has boundary value \(g_i(j)\) when restricted
to the smaller ball, so Dirichlet uniqueness makes the two extensions
agree there.  Their union is a smooth entire harmonic function \(U_i\).
At the sphere \(e^J\), its low projection has coordinates \(v_i(J)\).
Those columns are independent, so \(U_1,\ldots,U_p\) are independent.

\subsection{From the phase average to a bound at every point}

The graph component is orthogonal to the low component.  The column boxes
and \(\|V_j\|=O(\delta_j)\) therefore bound \(\|w_i(j)\|_2\) uniformly
above and away from zero.  Since the \(\alpha_i(j)\)'s lie in a fixed
compact positive interval,
\[
 \log\sigma_i(j)=\log\alpha_i(j)+O(\delta_j)
                =d_i(j)+O(\delta_j).
\]
It follows that
\begin{equation}
 \log\|g_i(j)\|_2
 =\sum_{\ell=J}^{j-1}d_i(\ell)
       +O\!\left(1+\sum_{\ell=J}^{j-1}\delta_\ell\right).
 \label{n:eq-mesh-log-growth}
\end{equation}
We evaluate the average on the right before passing to pointwise bounds.

For the frequency in \eqref{n:eq-continued-frequency}, define a periodic
primitive of its mean-zero part by
\[
 Q_i(s)=\int_0^s(\Theta_i(\sigma)-m_i)\,d\sigma.
\]
It is bounded because its change over \(P_{\mathrm{lab}}\) is zero.
Since \(s'(t)=t^{-3/4}\), integration by parts gives, for real \(T\geq J\),
\begin{align}
 \int_J^T(\Theta_i(s(t))-m_i)\,dt
 &=\bigl[t^{3/4}Q_i(s(t))\bigr]_J^T
       -\frac34\int_J^T t^{-1/4}Q_i(s(t))\,dt
 =O(T^{3/4}).
 \label{n:eq-phase-average}
\end{align}
The Lipschitz bound on \(\Theta_i\) makes the error between the integral
and its unit-mesh sum at most
\[
 C\sum_{j\leq T}\eta(j)=O(T^{1/4}).
\]
On the finite spectral range, \(b(j)\to a\) and smoothness of
\(\mathfrak d\) imply
\[
 d_i(j)-\Theta_i(s(j))\longrightarrow0
\]
uniformly in \(i\).  Its Cesaro mean consequently tends to zero.
Finally
\[
 \sum_{j\leq T}\delta_j
 =\sum_{j\leq T}j^{-9/16}=O(T^{7/16}).
\]
All four errors are sublinear in logarithmic time.  Combining them with
\eqref{n:eq-mesh-log-growth} gives
\begin{equation}
 \lim_{j\to\infty}\frac1j\log\|g_i(j)\|_2=m_i.
 \label{n:eq-mesh-exponent}
\end{equation}

There are only finitely many labels.  By \eqref{n:eq-subcritical-means},
choose one number
\[
 \max_{i\leq p}m_i<\gamma_*<k.
\]
After absorbing the finitely many initial steps into the constants,
\eqref{n:eq-mesh-exponent} implies
\begin{equation}
 \|g_i(j)\|_2\leq C_i e^{\gamma_* j}
 \qquad(j\geq J).
 \label{n:eq-mesh-bound}
\end{equation}
For any real \(t\in[j,j+1]\), \(U_i\) on the sphere \(e^t\) is the
restriction of the harmonic extension with outer trace \(g_i(j+1)\).
The all-radius contraction in Lemma~\ref{n:lem-inward-transfer} gives
\begin{equation}
 \|U_i(t,\cdot)\|_2
 \leq\|g_i(j+1)\|_2
 \leq C_i' e^{\gamma_* t}.
 \label{n:eq-all-radius-sphere-bound}
\end{equation}
The same contraction from \(e^t\) to \(e^j\) gives
\(\|g_i(j)\|_2\leq\|U_i(t,\cdot)\|_2\). Together the two bounds also
identify the exact spherical growth exponent:
\[
 \lim_{t\to\infty}\frac1t\log\|U_i(t,\cdot)\|_2=m_i.
\]
In particular, the upper estimate holds on every sufficiently large sphere.

The normalized divergence form of
\eqref{n:eq-logarithmic-harmonic} is uniformly elliptic on cylinders of
fixed width, with uniformly bounded coefficients in a finite sphere atlas.
The local \(L^2\)-to-sup estimate, applied to \(U_i\) and \(-U_i\) on
a cylinder of width four
\cite[Theorem~8.17, p.~194, with \(p=2\)]{GilbargTrudinger}, gives a
constant independent of large \(t\) such that
\[
 \sup_{\omega\in S^2}|U_i(t,\omega)|
 \leq C\left(\int_{t-2}^{t+2}
          \|U_i(\tau,\cdot)\|_2^2\,d\tau\right)^{1/2}
 \leq C_i''e^{\gamma_* t},
\]
where the last inequality uses \eqref{n:eq-all-radius-sphere-bound}.
The compact core is absorbed into the constant.
Proposition~\ref{n:prop-radial-realization} gives the exact distance identity
\(d_g(0,(r,\omega))=r=e^t\).  Since \(\gamma_*<k\), we conclude
\[
 |U_i(x)|\leq C_i'''(1+d_g(0,x))^k
 \qquad\text{for every }x\in\mathbb R^3.
\]
The single metric fixed at the product fixed point therefore has at least
\(p=(M+1)^2\) independent members of \(\mathcal H_k(\mathbb R^3,g)\).
Since \(p/(k+1)^2\to\vartheta^2>c\), this is at least \(c(k+1)^2\)
for all sufficiently large \(k\).  The construction was
made for the chosen sufficiently large \(k\); no common metric for all
degrees is asserted.


\begin{proof}[Completion of Theorem~\ref{n:main}]
Given \(4/9<v<1\) and \(1<c<9v/4\), set \(a=\sqrt v\), choose
\(\sqrt c<\vartheta<3a/2\), and then choose \(a^2<\kappa_*<1\).
Fix the equatorial bump and a sufficiently small convex near-round
neighborhood. It may be made arbitrarily small before \(k\) is chosen;
in particular, impose the strict frequency bound
\eqref{n:eq-near-round-frequency-bound}.
Lemma~\ref{n:lem-adjacent-doubles} and
Proposition~\ref{n:prop-angular-program} impose only cofinite conditions
on the integer \(k\). Choose one threshold above all of them and large
enough that \((\lfloor\vartheta k\rfloor+1)^2\geq c(k+1)^2\).
For each integer above this threshold choose its finite program, crossing
directions, gaps, and disjoint sign neighborhoods. Choose the Ricci
buffer \(C\) next. Proposition~\ref{n:prop-radial-realization} works for
every sufficiently large \(J\) after these choices. Enlarge \(J\) to meet
the freezing-slab, relative graph, box, and endpoint-sign requirements;
none imposes an upper bound on \(J\).

The simultaneous fixed point selects one control sequence for all \(p\)
lines, and the construction above gives \(p\geq c(k+1)^2\) independent
entire functions in \(\mathcal H_k\) for its metric \(g_k\). The radial
proposition gives smoothness, completeness, the Euclidean core,
nonnegative Ricci curvature, and positive Ricci curvature outside a
compact set. Section~\ref{n:geometric-consequences} gives
\(\operatorname{AVR}(g_k)=a^2=v\), uncountably many nonisometric pointed
cone limits, and the negative radial planes along a sequence tending to
infinity. These conclusions hold for every provisional control and hence
for the selected one. This proves the stated order of quantifiers, with
one threshold followed by each eligible \(k\) and then its metric \(g_k\).
\end{proof}

\begin{proof}[Proof of Corollary~\ref{n:near-euclidean}]
Fix \(\epsilon>0\) and \(1<c<9/4\). Choose
\[
 \max\{2\sqrt c/3,(1+\epsilon)^{-1}\}<a<1.
\]
Choose \(0<\delta<1\) small enough that
\[
 a^2(1-\delta)\geq(1+\epsilon)^{-2},
 \qquad 1+\delta\leq(1+\epsilon)^2.
\]
Run the preceding construction with \(v=a^2\), taking the initial
conformal neighborhood small enough for \eqref{n:eq-gauge-closeness}
with this \(\delta\). The reference program and its interpolation are
chosen with a positive margin inside that neighborhood, so all controls
remain there once \(J\) is sufficiently large. The global comparison
\eqref{n:eq-global-tensor-comparison} gives the asserted tensor bounds.
Integrating along curves and taking infima gives the stated distance
bounds. The choice of \(a\) can be arbitrarily close to \(1\), so the
volume ratio can be prescribed arbitrarily close to \(1\) as well.
\end{proof}
